\section{Integral form of Heaviside theta}
\label{appendix:heaviside}

We're going to show that,
\begin{align}
    \theta\qty(x) = -\frac{1}{2\pi\im}\lim\limits_{z\rightarrow0^+}\int\limits_{-\infty}^{+\infty}\dd{z}\frac{\e^{-\im x z}}{z+\im\epsilon},\ \ \ x\in\mathbb R^*
\end{align}

But first, let's take a look at the following function,
\begin{align}
    f\qty(x)=\int\limits_{-\infty}^{x}\dd{y}\delta\qty(y),\ \ \ x\in\mathbb R^*
\end{align}
if $x>0$ then the integration domain passes through zero, and by the Dirac delta definition the integral gives one, if $x<0$ the integration domain 
doesn't contains zero, hence the integral gives zero. This is the Heaviside theta function!
\begin{align}
   \theta\qty(x)=\int\limits_{-\infty}^{x}\dd{y}\delta\qty(y)=\begin{cases}1,&x>0\\0,&x<0\end{cases}
\end{align}
from this representation we can work out a little bit more by inserting another representation of the Dirac delta. Sadly, because we aren't dealing with usual converging integrals, 
the Dirac delta isn't a function, isn't finite and because it's integral representation does not converge in the usual sense we'll have to add a term which helps improve convergence, 
but does not changes the value of the function,
\begin{align}
   \theta\qty(x)&=\int\limits_{-\infty}^{x}\dd{y}\e^{\epsilon y}\delta\qty(y)
\end{align}
any value of $\epsilon$ is a good choice, as when integrated with the Dirac delta it returns just $\e^{\epsilon 0} =1 $. But as our integral goes all the way to $y\rightarrow-\infty$, 
it's more interesting to keep $\epsilon>0$. As the final result is just independent of this additional variable, we can as well consider it as a limit,
\begin{align}
   \theta\qty(x)&=\lim\limits_{\epsilon\rightarrow0^+}\int\limits_{-\infty}^{x}\dd{y}\e^{\epsilon y}\delta\qty(y)
\end{align}
despite it doing nothing by now it will help improve the convergence as soon as we substitute the integral representation of the Dirac delta,
\begin{align}
   \theta\qty(x)&=\lim\limits_{\epsilon\rightarrow0^+}\int\limits_{-\infty}^{x}\dd{y}\e^{\epsilon y}\frac{1}{2\pi}\int\limits_{-\infty}^{+\infty}\dd{z}\e^{-\im y z}\\
   \theta\qty(x)&=\lim\limits_{\epsilon\rightarrow0^+}\int\limits_{-\infty}^{x}\dd{y}\frac{1}{2\pi}\int\limits_{-\infty}^{+\infty}\dd{z}\e^{-\im y\qty( z+\im\epsilon)}
\end{align}
now we recklessly exchange the order of integration and perform the $y$ integration,
\begin{align}
   \theta\qty(x)&=\lim\limits_{\epsilon\rightarrow0^+}\frac{1}{2\pi}\int\limits_{-\infty}^{+\infty}\dd{z}\frac{\e^{-\im y\qty( z+\im\epsilon)}}{\qty(-\im)\qty(z+\im\epsilon)}\eval_{y=-\infty}^{y=x}\\
   \theta\qty(x)&=-\lim\limits_{\epsilon\rightarrow0^+}\frac{1}{2\pi \im}\int\limits_{-\infty}^{+\infty}\dd{z}\frac{1}{{z+\im\epsilon}}\qty(\e^{-\im x\qty( z+\im\epsilon)}-\lim\limits_{y\rightarrow -\infty}\e^{-\im y\qty( z+\im\epsilon)})
\end{align}
the limit inside the integral is zero, because $\epsilon>0$, then,
\begin{align}
   \theta\qty(x)&=-\lim\limits_{\epsilon\rightarrow0^+}\frac{1}{2\pi \im}\int\limits_{-\infty}^{+\infty}\dd{z}\frac{\e^{-\im x\qty( z+\im\epsilon)}}{{z+\im\epsilon}}
\end{align}
Notice that now the $\epsilon$ limit is in fact essential, it does not commute anymore with the integral, as the integrand has a pole in this limit. Despite this, we can neglect the $\epsilon$ 
present in the exponential, as it does not introduces new poles neither cross poles or branch cuts. The final expression we arrive is then,
\begin{align}
   \theta\qty(x)&=-\lim\limits_{\epsilon\rightarrow0^+}\frac{1}{2\pi \im}\int\limits_{-\infty}^{+\infty}\dd{z}\frac{\e^{-\im xz}}{{z+\im\epsilon}}
\end{align}
as desired.